% AssemblyP1 white paper.
% Reproducible build: see paper/README.md and paper/build.sh
\documentclass[11pt,a4paper]{article}

% Shared preamble for the AssemblyP1 white paper.
% Conservative package set; see paper/manifest.scm.

\usepackage[T1]{fontenc}
\usepackage[utf8]{inputenc}
\usepackage{lmodern}
\usepackage[margin=1in]{geometry}
\usepackage{amsmath}
\usepackage{amssymb}
\usepackage{amsthm}
\usepackage{microtype}
\usepackage{booktabs}
\usepackage{enumitem}
\usepackage{xcolor}
\usepackage{tikz}
\usetikzlibrary{arrows.meta,positioning,fit,backgrounds}
\usepackage[backend=biber,style=authoryear,maxbibnames=99]{biblatex}
\usepackage[hidelinks,colorlinks=true,linkcolor=blue!45!black,citecolor=green!40!black,urlcolor=blue!55!black]{hyperref}

% ---------- Theorem environments ----------
\theoremstyle{plain}
\newtheorem{theorem}{Theorem}[section]
\newtheorem{proposition}[theorem]{Proposition}
\newtheorem{lemma}[theorem]{Lemma}
\newtheorem{corollary}[theorem]{Corollary}

\theoremstyle{definition}
\newtheorem{definition}[theorem]{Definition}
\newtheorem{example}[theorem]{Example}
\newtheorem{counterexample}[theorem]{Counterexample}
\newtheorem{conjecture}[theorem]{Conjecture}
\newtheorem{problem}[theorem]{Problem}

\theoremstyle{remark}
\newtheorem{remark}[theorem]{Remark}

% ---------- Epistemic-status markup ----------
% Every nontrivial claim carries one of these labels, inline or after a proof.
\newcommand{\statustag}[1]{\textnormal{\textsc{[#1]}}}
\newcommand{\srcfact}{\statustag{source fact}}
\newcommand{\srcinf}{\statustag{source-supported inference}}
\newcommand{\projchoice}{\statustag{project modeling choice}}
\newcommand{\mathfact}{\statustag{mathematical result}}
\newcommand{\kernelcheck}{\statustag{kernel-checked}}
\newcommand{\computed}{\statustag{verified computation}}
\newcommand{\boundedev}{\statustag{bounded computational evidence}}
\newcommand{\openstatus}{\statustag{open}}
\newcommand{\srcgap}{\statustag{source gap}}

\newcommand{\statusline}[1]{\par\smallskip\noindent\textit{Status.}~#1}
\newcommand{\provenance}[1]{\par\smallskip\noindent\textit{Provenance.}~#1}

% A small legend used in several places.
\newcommand{\statuslegend}{%
  \srcfact, \srcinf, \projchoice, \mathfact, \kernelcheck, \computed,
  \boundedev, \srcgap, and \openstatus.}

% ---------- Notation ----------
\newcommand{\Alphabet}{\Sigma}
\newcommand{\spec}{\operatorname{spec}}
\newcommand{\supp}{\operatorname{supp}}
\newcommand{\occ}{\operatorname{occ}}
\newcommand{\KL}{\operatorname{KL}}
\newcommand{\ellpop}{\ell}
\newcommand{\reals}{\mathbb{R}}
\newcommand{\nats}{\mathbb{N}}
\newcommand{\ints}{\mathbb{Z}}
\newcommand{\Is}{I_s}
\newcommand{\rc}{\operatorname{rc}}
\newcommand{\Lik}{\mathcal{L}}
\newcommand{\lendist}{\mathcal{D}}
\newcommand{\cands}{\mathcal{C}}

% Semantic emphasis for the two conclusion schemas.
\newcommand{\weak}{\mathrm{W}}   % truth is a maximizer
\newcommand{\strong}{\mathrm{S}} % every maximizer is the truth up to equivalence

% ---------- Lightweight figure styles ----------
\tikzset{
  genomebase/.style={circle,draw,inner sep=1pt,minimum size=1.35em,font=\small},
  readarc/.style={-{Latex[length=2mm]},thick},
}


\addbibresource{references.bib}

\title{\bfseries Bridging Conditions and Maximum-Likelihood\\
Genome Assembly\\[0.4em]
\large A Source-Faithful Study of a 2016 Open Question}

\author{AssemblyP1 Project\thanks{This manuscript is a research artifact of the
\href{https://github.com/ottojung/assemblyp1}{assemblyp1} repository. It
distinguishes source facts, project interpretations, proved theorems, and
kernel-checked finite certificates throughout; see
Appendix~\ref{app:ledger} for the consolidated epistemic ledger.}}

\date{\today}

\begin{document}

\maketitle

\begin{abstract}
\noindent
Shomorony, Kim, Courtade, and Tse (2016) asked whether the repeat-bridging
conditions that make shotgun assembly \emph{information-feasible} also guarantee
that ``the maximum-likelihood sequence is the true sequence,'' referring to the
maximum-likelihood assembly formulation of Medvedev and Brudno (2009). We give a
source-faithful mathematical study of this question that keeps four layers
visibly separate: the historical definitions, the ambiguities those definitions
leave open, theorems and counterexamples about the resulting formulations, and
new models introduced by this project to repair the weaknesses found.

We show that the published sentence does not by itself determine several
modeling axes --- which Medvedev--Brudno likelihood layer is meant, whether reads
are oriented or reverse-complement molecules, whether candidates are
length-restricted, and whether the conclusion asserts existence of a maximizer
or uniqueness of the maximizing genome. Along these axes we record two kinds of
concrete results. On the negative side, kernel-checked same-length witnesses
refute the maximizer conclusion under the exact multinomial and under the
fixed-length binomial approximation, over all circular candidates, and further
witnesses refute it under the reverse-complement molecule reading of the
Medvedev--Brudno Section~6.2 flow. On the positive side, we prove a rigidity
theorem: under Shomorony's information-feasibility condition and strict oriented
read types, the truth's length-$L$ spectrum is the unique positive circulation
of its length-$L$ support graph, so no strict same-length Section~6.2
counterexample exists for any genome length, read length, or alphabet; on that
truth-admissible row the same complete-spectrum input used below makes the truth
unique up to cyclic rotation.

We then present a \emph{new}, project-level model that isolates the statistical
failure mechanism and removes it. Under population (infinite-read) maximum
likelihood, the truth is always a maximizer by a Gibbs/KL identity, so
uniqueness becomes pure spectrum identifiability. For the project's oriented,
primitive, P2-admissible candidate class, a reduction to the
complete-spectrum uniqueness theorem of Bresler, Bresler, and Tse (2013)
yields uniqueness up to cyclic rotation. This repaired result does not settle the
2016 finite-sample question, and we say exactly why.
\end{abstract}

\tableofcontents

\section{Introduction}
\label{sec:intro}

\subsection{The question}

Modern genome assembly reconstructs an unknown genome from many short
overlapping reads. Two lines of theory ask two different questions about this
problem. \emph{Information-feasibility} asks whether the reads contain enough
information to determine the genome; the answer depends on the repeat structure
of the genome and on the read length, and Shomorony, Kim, Courtade, and Tse
\cite{shomorony2016} gave repeat-bridging conditions under which a specific
algorithm recovers the true circular sequence up to cyclic shift. \emph{Maximum
likelihood} instead asks which genome is the most probable source of the
observed reads; Medvedev and Brudno \cite{medvedev2009} formulated this as an
optimization problem over candidate genomes and solved a relaxed version by
min-cost flow.

Shomorony et al.\ close their Discussion with a question that sits between the
two lines \cite[Discussion]{shomorony2016}:

\begin{quote}
``Understanding whether bridging conditions can be used to guarantee that the
maximum-likelihood sequence is the true sequence is currently an open
question.''
\end{quote}

This manuscript studies that question as a mathematical problem --- but not by
pretending it is a single well-posed theorem. The sentence names a likelihood
formulation only by citation and leaves several choices open. Our first
contribution is to make those choices explicit; our second is to report what is
provable or refutable for each resulting formulation.

\subsection{Why maximum likelihood is not automatic}

A naive argument says the question should be trivial. The reads were generated
from the true genome, so the true genome should maximize their likelihood. The
argument fails for a structural reason: the read distribution can be
\emph{non-identifiable}. The reads are drawn from the multiset of length-$L$
windows of the genome, so two genomes with the same window spectrum induce
exactly the same read distribution and always tie.

\begin{example}[Structural ambiguity before the repeat boundary]
\label{ex:ABAB}
Let the alphabet be $\{A,B\}$ and let $L=2$. The circular genomes
\[
  S = AAABAB, \qquad D = AABAAB
\]
of length $6$ are not cyclic rotations of one another, yet their circular
$2$-mer spectra agree:
\[
  \spec_2(S) = \spec_2(D) = \{\,AA:2,\ AB:2,\ BA:2\,\}.
\]
Indeed $S$ has windows $AA,AA,AB,BA,AB,BA$ and $D$ has windows
$AA,AB,BA,AA,AB,BA$. Under uniform-start error-free sampling the observed read
distribution is the same for both genomes, so \emph{no} likelihood objective
--- exact, approximate, or flow-based --- can prefer $S$ to $D$ here.
\statusline{\srcfact{} for the spectra and the rotation claim;
\mathfact{} for the resulting tie. This is the ``$AAABAB/AABAAB$''
example recorded in the project synthesis.}
\end{example}

The lesson of Example~\ref{ex:ABAB} is that maximum likelihood cannot manufacture
identifiability. Bridging conditions exist precisely to remove the structural
ambiguity: they are hypotheses on the repeat structure and read coverage of the
\emph{true} genome that, in Shomorony et al.'s theorem, force the true genome to
be the unique Eulerian object consistent with the reads. The published question
asks whether those same hypotheses also make the true genome optimal for a
likelihood objective. It is therefore a question about whether a
\emph{reconstruction} guarantee transfers to an \emph{optimization} guarantee.

\subsection{What the question does not say}

Before any theorem can be attributed to the published sentence, four modeling
axes must be fixed, and we will argue in Section~\ref{sec:open} that the
sentence fixes none of them conclusively. The named formulation of
\cite{medvedev2009} contains at least four distinct objects:

\begin{enumerate}[label=(\roman*)]
  \item the exact global read-count multinomial with candidate-intrinsic length;
  \item its separable product-of-binomial-marginals approximation, which replaces
        the candidate length by an external genome size;
  \item the Section~6.2 bidirected read-overlap-graph flow, whose feasible
        objects are flows and need not be sequences; and
  \item the underlying broad ``maximum-likelihood principle.''
\end{enumerate}

Independently, the sentence does not say whether reads are oriented strings or
reverse-complement molecule classes; whether candidate genomes are restricted to
the true length; or whether the conclusion asserts that the truth is \emph{some}
maximizer or that \emph{every} maximizer is the truth up to a genome equivalence.
These axes matter, and we treat them as part of the mathematical object rather
than as notation.

\subsection{Contributions}

This manuscript does four things.

\paragraph{Negative finite results.}
We exhibit kernel-checked finite counterexamples showing that the
information-feasibility hypothesis does \emph{not} force the truth to be a
likelihood maximizer under several precise readings of the 2016 question: the
exact multinomial over all circular candidates (Theorem~\ref{thm:cex-exact}), the
fixed-length binomial approximation over all circular candidates
(Theorem~\ref{thm:cex-binomial}), and the reverse-complement molecule reading of
the Section~6.2 flow (Theorem~\ref{thm:cex-molecule}). A candidate-set inclusion
lemma (Lemma~\ref{lem:negative-transfer}) shows that a same-length refutation
transfers to every superclass of candidates, so these results are not confined
to length-restricted comparisons.

\paragraph{A positive same-length rigidity theorem.}
Under \emph{strict oriented} read types, Shomorony's $\Is$, and the
same-length spelled-candidate restriction, the truth's spectrum is the unique
positive circulation of its window-support graph (Theorem~\ref{thm:rigidity}).
Consequently every same-length Section~6.2 spelled candidate ties the truth under
any spectrum-based objective, and no strict same-length counterexample exists for
\emph{any} genome length, read length, or alphabet. The equal-length restriction
is essential: allowing one extra symbol reintroduces a strict witness
(Example~\ref{ex:variable-length}).

\paragraph{A repaired population model.}
We present a new model, motivated by the finite failures rather than by the
historical papers: replace empirical read counts by the exact population read
distribution. A Gibbs/KL identity (Lemma~\ref{lem:gibbs}) makes the truth a
population maximizer \emph{unconditionally}. Uniqueness then reduces to spectrum
identifiability, and for the project's oriented, primitive, P2-admissible
candidate class this is supplied by a complete-spectrum uniqueness theorem of
Bresler, Bresler, and Tse \cite{bresler2013}. The resulting
Theorem~\ref{thm:population} is the current positive high point of the project.
It does not settle the finite 2016 question, and we state the gap explicitly.

\paragraph{Exposition as a first-class concern.}
Because the mathematical content is an interface between under-specified
sources and project-level repairs, correctness includes keeping the epistemic
layers separate. Theorem statements name their assumption surface; every
nontrivial claim carries a status tag (\statuslegend); and
Appendices~\ref{app:modeling} and~\ref{app:ledger} explain the modeling choices
and consolidate the status of each claim.

\subsection{How to read this paper}

Section~\ref{sec:model} fixes notation and the source-faithful sequencing and
bridging model. Section~\ref{sec:open} states the published question and
determines what the sources do and do not fix. Section~\ref{sec:finite} reports
the results about the literature-derived finite formulations.
Section~\ref{sec:population} introduces the population model and its uniqueness
theorem. Section~\ref{sec:discussion} draws the two failure modes together and
lists what remains open. Appendix~\ref{app:modeling} explains the modeling
choices deliberately; Appendix~\ref{app:ledger} is the consolidated epistemic
ledger; Appendix~\ref{app:reproduction} records how to rebuild the paper and
re-run every computational certificate.

A reader who wants only the mathematical story can read
Sections~\ref{sec:model}--\ref{sec:discussion} and treat the appendices as
reference. A reader who wants to scrutinize the source correspondence should
read Appendix~\ref{app:modeling} first.

\section{The sequencing and bridging model}
\label{sec:model}

This section fixes notation and records the source-faithful model. We separate
the parts that come from the literature from the project's normalizations; the
former are tagged \srcfact{} and the latter \projchoice. The de Bruijn/Eulerian
view of assembly that underlies both the repeat theory and the Section~6.2 graph
is classical \cite{idury1995,pevzner2001}. A fuller discussion of
\emph{why} the model has this shape is deferred to Appendix~\ref{app:modeling}.

\subsection{Circular genomes and spectra}

\begin{definition}[Circular word, window, spectrum]
\label{def:spectrum}
Fix a finite alphabet $\Alphabet$ of size $\sigma$. A \emph{circular word} $D$ of
length $n=|D|$ is a function $D:\ints/n\ints \to \Alphabet$, written
$D_0D_1\cdots D_{n-1}$ with indices mod $n$. For $L \le n$ the \emph{length-$L$
window} at start $t$ is
\[
  W_t(D) \;=\; D_t D_{t+1}\cdots D_{t+L-1},
\]
read cyclically. The \emph{length-$L$ spectrum} of $D$ is the function
\[
  \spec_L(D)(w) \;=\; \#\{\, t \in \ints/n\ints : W_t(D)=w \,\}
\]
on length-$L$ words, with total $\sum_w \spec_L(D)(w) = n$. Its support is
$\supp\spec_L(D) = \{\,w : \spec_L(D)(w) > 0\,\}$.
\statusline{\srcfact{} for the circular, oriented-window model
\cite[\S2]{shomorony2016}; the notation is ours.}
\end{definition}

Two circular words are \emph{cyclically equivalent}, written $D \sim D'$, if one
is a rotation of the other. A circular word is \emph{primitive} if its minimal
period equals its length, i.e.\ it is not $P^k$ for a shorter word $P$ and
$k\ge 2$. Cyclic equivalence preserves $|D|$ and $\spec_L(D)$, so it is the
weakest equivalence that a likelihood on spectra can ever distinguish.

\begin{remark}[Why cyclic shift is unavoidable]
\label{rem:cyclic}
Rotation is a bijection of start positions commuting with $t \mapsto W_t(D)$;
hence $\spec_L(\rho D)=\spec_L(D)$ for every rotation $\rho D$ and every $L$.
A circular candidate therefore has (at least) $n$ spellings with identical
spectra, and a conclusion that ignores cyclic shift can never be true for a
non-rotation-fixed genome. This is a \mathfact, not a modeling preference.
\end{remark}

\subsection{Reads and the information-feasible hypothesis}

\begin{definition}[Read realization and coverage]
\label{def:reads}
A \emph{read realization} of the circular truth $S$ with read length $L$ is a
multiset $R$ of start positions, possibly with repetition; the observed read
strings are the corresponding windows $W_t(S)$. It \emph{covers} $S$ if every
position of $S$ lies in at least one read interval. The observable data are the
read \emph{multiplicities} $x(w) = \#\{t \in R : W_t(S)=w\}$, not merely the set
of distinct read strings.
\statusline{\srcfact{} for uniform independent error-free sampling and coverage
\cite[\S2]{shomorony2016}; retaining multiplicity is required for likelihood
evaluation.}
\end{definition}

We now record the repeat and bridging vocabulary. The source-faithful reading is
that Shomorony et al.\ inherit these definitions from Bresler, Bresler, and Tse
\cite{bresler2013}; we follow that attribution.

\begin{definition}[Repeat objects with maximality]
\label{def:repeats}
A length-$\ell$ \emph{repeat pair} consists of two distinct start positions whose
length-$\ell$ windows are equal and which are \emph{maximal} on both sides: the
symbols immediately preceding the two copies differ and the symbols immediately
following them differ. A \emph{triple repeat} uses three starts with equal
length-$\ell$ windows, with the three-copy maximality condition that the
preceding symbols are not all equal and the following symbols are not all equal.
A repeat pair is \emph{interleaved} if its four selected starts alternate in
cyclic order; its length is the shorter constituent length.
\statusline{\srcfact{} \cite[repeat definitions]{bresler2013};
the cyclic-order normalization of interleaving is a \projchoice{} proved
equivalent to the source's displayed linear orders on a circle.}
\end{definition}

A read interval $[r,r+L)$ \emph{bridges} a lifted copy $[t,t+\ell)$ iff
$r<t$ and $t+\ell<r+L$, i.e.\ the read strictly extends beyond the copy on both
sides. A repeat is \emph{bridged} if at least one selected copy is bridged; a
triple repeat is \emph{all-bridged} if \emph{every} selected copy is bridged.
This must not be confused with the weaker ``at least one copy bridged'' reading
of a triple repeat.

\begin{definition}[Information-feasible set $\Is$]
\label{def:Is}
A read realization $R$ of $S$ is \emph{information-feasible}, written $R \in
\Is$, if
\begin{enumerate}[label=(C\arabic*),leftmargin=*]
  \item $R$ covers $S$;
  \item every triple repeat of $S$ is all-bridged;
  \item every interleaved repeat pair of $S$ is bridged.
\end{enumerate}
\statusline{\srcfact{} \cite[Eq.~(1)]{shomorony2016}, inheriting
\cite{bresler2013}. The triplet is the hypothesis side of the published
open-question sentence.}
\end{definition}

A length-$L$ read can bridge a copy of a repeat of length $\ell$ only if
$\ell \le L-2$: strict extension on both sides requires the repeat to fit
strictly inside the read. We use this repeatedly.

\subsection{Observation models and candidate classes}

The likelihood layer must consume only observable read multiplicities plus a
candidate genome; it must not see the latent true placements. We keep the
following three objectives distinct, because the published sentence does not
select among them (Section~\ref{sec:open}).

\begin{definition}[Exact multinomial; binomial approximation]
\label{def:likelihoods}
Let $x$ be observed read-type multiplicities with $n=\sum_w x(w)$, and let $D$ be
a circular candidate.
\begin{enumerate}[label=(\roman*),leftmargin=*]
  \item The \emph{exact multinomial} objective is
  \[
    \Lik_{\mathrm{ex}}(D;x)
      = \frac{n!}{\prod_w x(w)!}\,
        \prod_{w\,:\,\spec_L(D)(w)>0}
        \left(\frac{\spec_L(D)(w)}{|D|}\right)^{x(w)},
  \]
  with $|D|$ the candidate-intrinsic length, and value $0$ if some observed
  word has zero multiplicity in $D$.
  \item The \emph{fixed-length binomial} objective with external genome size $N$
  is
  \[
    \Lik_{\mathrm{bin}}(D;x)
      = \prod_w \binom{n}{x(w)}
        \left(\frac{\spec_L(D)(w)}{N}\right)^{x(w)}
        \left(1-\frac{\spec_L(D)(w)}{N}\right)^{n-x(w)},
  \]
  defined when $0 \le \spec_L(D)(w) \le N$ for all $w$.
\end{enumerate}
\statusline{(i) is the \srcfact{} exact global read-count likelihood
\cite[\S6.1]{medvedev2009}; (ii) is the \srcfact{} separable approximation that
replaces the candidate length $N(D)$ by an external, known $N$, which Medvedev
and Brudno explicitly adopt for their algorithm. The two are not
definitionally equal. We refer to (ii) as the \emph{Section~6.1 product} and
write $\Lik_{6.1} = \Lik_{\mathrm{bin}}$; the formula is the same when the
read-type space is taken to be molecule classes rather than oriented windows.}
\end{definition}

A \emph{candidate class} is a set $\cands$ of circular words to which the
objective is restricted. The following classes occur below:
\begin{itemize}[leftmargin=*]
  \item all nonempty circular words (where the exact objective is defined);
  \item the same-length class $\{\,D : |D|=|S|\,\}$;
  \item Section~6.2 \emph{spelled candidates}: circular words $D$ whose window
        support equals the observed support, $\supp\spec_L(D)=\supp x$, reflecting
        the source's per-vertex lower bound of $1$ on each read molecule
        \cite[\S6.2, Observation 7]{medvedev2009};
  \item the project-level classes P1 and P2 defined in
        Section~\ref{sec:population}.
\end{itemize}

Following the repository's formalization contract, we never silently identify
these classes. A theorem proved over a restricted class is a restricted result;
but a \emph{negative} result over a subclass transfers to every superclass
(Lemma~\ref{lem:negative-transfer}).

\subsection{The two conclusion schemas}

The English phrase ``the maximum-likelihood sequence is the true sequence'' has
two inequivalent formalizations, and the sources state no tie rule.

\begin{definition}[Maximizer and uniqueness schemas]
\label{def:schemas}
Fix an objective $\Lik$, an observation $x$ generated from truth $S$, a
candidate class $\cands$ containing $S$, and a genome equivalence $\approx$.
\begin{align*}
  \weak &: \quad \forall D \in \cands,\quad \Lik(D;x) \le \Lik(S;x);\\
  \strong &: \quad \weak \ \wedge\ \forall D \in \cands,\quad
             \Lik(D;x)=\Lik(S;x) \Rightarrow D \approx S.
\end{align*}
The weak schema asserts that the truth is \emph{a} maximizer; the strong schema
adds that \emph{every} maximizer is the truth up to $\approx$.
\end{definition}

Two facts about the equivalence are forced rather than chosen.
Cyclic shift must belong to $\approx$ for the strong schema to be satisfiable at
all (Remark~\ref{rem:cyclic}). And reverse-complement equivalence is forced
\emph{exactly} when read types are reverse-complement molecule classes: under
molecule types a candidate and its reverse complement always tie, while under
oriented read types they need not. The read-type convention and the genome
equivalence are therefore one coupled choice, not two independent parameters.
We return to this in Section~\ref{sec:open}.

Finally, a strict inequality makes equivalence and tie conventions irrelevant: if
some candidate has $\Lik(D;x) > \Lik(S;x)$ strictly, then both $\weak$ and
$\strong$ are false under \emph{every} equivalence and tie convention. The
kernel-checked witnesses of Section~\ref{sec:finite} are of this strict kind,
which is why they are robust to the unresolved equivalence choices.

\section{The published question and what its sources do not say}
\label{sec:open}

\subsection{The exact sentence and its referent}

The target sentence is \cite[Discussion, p.~i501]{shomorony2016}:

\begin{quote}
``The maximum-likelihood formulation of the AP (Medvedev and Brudno, 2009), on
the contrary, seems to be robust to these issues, and thus a good candidate for
the `correct' formulation. Understanding whether bridging conditions can be used
to guarantee that the maximum-likelihood sequence is the true sequence is
currently an open question.''
\end{quote}

A full-text scan of the accepted article finds no section, equation, formula, or
candidate-class pointer into Medvedev and Brudno. The sentence names the work
only by bibliography and the phrase ``maximum-likelihood formulation.'' The
surrounding paragraph contrasts information-feasible reconstruction with
``optimization-based formulation[s] of the AP such as those considered by
Nagarajan and Pop (2009) \dots\ and Medvedev and Brudno (2009)''
\cite{nagarajan2009,medvedev2009}; it fixes no objective and no tie rule. This is
the central source fact of this section: \srcfact{} the 2016 text does not
select a referent, a read-type convention, a candidate length, or a
maximizer/uniqueness reading.

\subsection{Four Medvedev--Brudno objects}

The cited paper contains four materially different objects, and they answer the
word ``sequence'' differently.

\begin{table}[t]
\centering
\small
\begin{tabular}{@{}p{0.30\textwidth}p{0.42\textwidth}p{0.20\textwidth}@{}}
\toprule
Reading & What the optimization returns & Sequence? \\
\midrule
(1) exact multinomial, candidate-intrinsic $N(D)$
  & a circular genome $D$
  & yes \\
(2) fixed-$N$ binomial approximation
  & a circular genome $D$ scored through per-type counts
  & yes \\
(3) Section~6.2 bidirected flow
  & a flow / ``(non-contiguous) assembly''
  & not necessarily \\
(4) broad ML principle
  & unspecified
  & n/a \\
\bottomrule
\end{tabular}
\caption{The four referents of ``the maximum-likelihood formulation
(Medvedev and Brudno, 2009).'' Reading (1) is what the paper \emph{names} as its
global read-count likelihood but explicitly abandons as not separable;
reading (2) is what its algorithm \emph{does} and what the contemporaneous
literature treats as the method; reading (3) is its algorithm's feasible set;
reading (4) is what the 2016 sentence most safely says literally. No primary
source selects one (\srcfact); the best operational sequence-level fit is (2)
(\srcinf).}
\label{tab:referents}
\end{table}

Reading (3) deserves emphasis: its feasible objects are flows, and Medvedev and
Brudno route the question of a single assembled sequence to a later heuristic.
Thus if the 2016 phrase meant the Section~6.2 flow literally, ``the
maximum-likelihood \emph{sequence}'' would be a category error. That is evidence
against reading (3) as the direct referent, while not excluding it as a
formalization of the underlying optimization.

Reading (2) is not a project invention. The dissertation version of the work
presents the fixed-size method as the operative one \cite{medvedev2010}, a
contemporaneous survey describes the Medvedev--Brudno assembler as requiring the
accurate genome size as a parameter \cite{howison2013}, and a follow-up
formulation advertises its improvement precisely as better genome-size handling
\cite{varma2011}. This is evidence about how the community read the method; it
does not upgrade reading (2) to the literal referent of the 2016 sentence.

A refinement matters for negative results. The binomial objective contains
$\log d_i$ and $\log(N-d_i)$, so it is defined only when $0<d_i<N$ for every
type; the exact multinomial is defined on every nonempty candidate. A refutation
therefore has to name the objective's domain, not just the candidate class.

\subsection{The axes the sentence leaves open}

\begin{table}[t]
\centering
\small
\begin{tabular}{@{}p{0.30\textwidth}p{0.36\textwidth}p{0.26\textwidth}@{}}
\toprule
Axis & Determination & Status \\
\midrule
named likelihood layer & readings (1)--(4) all literal & \srcgap \\
read-type space & oriented windows vs.\ molecule classes & \srcgap \\
candidate length & not fixed by the sentence; fixed length is best
  operational fit & \srcgap \\
maximizer vs.\ uniqueness & both literal & \srcgap \\
cyclic shift in $\approx$ & \textbf{forced} for the strong schema & \mathfact \\
reverse complement in $\approx$ & \textbf{forced iff} molecule read types
  & \mathfact\ + \srcgap \\
other genome equivalences & no source support & \srcfact \\
tie rule & none stated; many optima possible & \srcfact \\
\bottomrule
\end{tabular}
\caption{What the published sentence does and does not determine. The two
``forced'' rows are mathematical consequences rather than readings; the
remaining axes are genuine source gaps.}
\label{tab:axes}
\end{table}

Three points from Table~\ref{tab:axes} are worth isolating.

\paragraph{Candidate length is not the same as known genome size.}
Medvedev and Brudno's Section~6.1 approximation supplies an external genome-size
parameter $N$ to the likelihood. It does \emph{not} thereby restrict every
candidate $D$ to $|D|=N$. Knowing or supplying $N=|S|$ and restricting
candidates to $|D|=|S|$ are different assumptions; a same-length theorem is a
restricted result unless a source argument fixes the restriction.

\paragraph{Cyclic shift is mandatory, reverse complement is coupled.}
By Remark~\ref{rem:cyclic} the strong schema requires cyclic shift in $\approx$.
Reverse complement is forced precisely when read types are molecule classes: a
molecule and its reverse complement are always tied under class-indexed
observations, while under oriented reads the two count vectors generally differ.
Hence ``oriented with cyclic shift'' and ``molecular with dihedral equivalence''
are the two coherent source-supported panels; they are not freely mixable.

\paragraph{Ties are generic.}
The exact circular likelihood is invariant under cyclic shift, so distinct
spellings tie. A contemporaneous formulation states that many Eulerian tours can
have equal likelihood, so ``the maximum-likelihood sequence'' cannot be assumed
unique without an equivalence convention.

\subsection{Residual source gaps}

We record two gaps that bear on how strongly any formal result can be attributed
to the published sentence. First, the accepted publisher supplement of
\cite{shomorony2016} --- the last unexamined accepted artifact that could contain
a likelihood definition --- has not been retrieved through the available
execution path; a definition located there could narrow the referent. Second,
Medvedev and Brudno's Section~6.1 writes that there are ``$4^k$ such variables''
while its model vocabulary is $k$-molecules (reverse-complement classes), so the
index set and the vocabulary are in internal tension. A formalization must name
which it uses. Neither gap changes the mathematics below; both change how
directly a result may be advertised as settling the historical question.

\paragraph{Consequence for this paper.}
Because no referent is selected, we do not present any single theorem as ``the''
answer to the 2016 question. Instead, Section~\ref{sec:finite} proves and refutes
statements for named points in the axis space, and Section~\ref{sec:population}
presents a project-level repaired model that is explicitly not a historical
claim. A formal statement counts as a settlement of the published problem only
after an argument fixes the axes of Table~\ref{tab:axes} to a source-supported
panel; that argument is not currently available.

\section{Finite results for the literature-derived formulations}
\label{sec:finite}

This section reports what is known about the question when the observation model
is a finite sample and the objective is one of the Medvedev--Brudno objects of
Section~\ref{sec:model}. The results are of two kinds. First a \emph{positive}
rigidity theorem for same-length oriented Section~6.2 candidates under
Shomorony's $\Is$. Then a sequence of \emph{negative} witnesses, each refuting a
named, precisely scoped statement. We begin with the logical device that makes
the negatives robust.

\subsection{Negative results transfer; positive results do not}

\begin{lemma}[Negative transfer]
\label{lem:negative-transfer}
Let $\Lik(\cdot;x)$ be an objective defined on a candidate class $\cands$, let
$\cands'\subseteq\cands$, and let $S,D\in\cands'$. If
$\Lik(D;x) > \Lik(S;x)$ then $S$ is not a maximizer of $\Lik(\cdot;x)$ over
$\cands$.
\end{lemma}

\begin{proof}
$D\in\cands$ is a candidate with strictly larger value than $S$, so the maximum
over $\cands$ exceeds $\Lik(S;x)$.
\end{proof}

The converse fails: a positive result over a restricted class does not imply the
same over a superclass, because the superclass may contain a better candidate.
This asymmetry is why a same-length \emph{negative} witness that is
candidate-set closed settles the unrestricted question negatively, while a
same-length \emph{positive} theorem is a restricted result unless the restriction
is itself source-justified.

\subsection{A positive same-length oriented rigidity theorem}

We work under the strict oriented single-strand convention: read types are the
oriented length-$L$ windows of $S$ with no reverse-complement collapse, and a
same-length Section~6.2 spelled candidate is a circular word $D$ of length
$G=|S|$ with $\supp\spec_L(D)=\supp\spec_L(S)$.

\begin{definition}[Window-support graph]
\label{def:graph}
Given a circular word $S$ and read length $L$, the graph $X_S$ has one node for
each distinct length-$(L-1)$ window of $S$, and one directed edge for each
distinct length-$L$ window $w\in\supp\spec_L(S)$, running from
$\operatorname{prefix}_{L-1}(w)$ to $\operatorname{suffix}_{L-1}(w)$ and carrying
multiplicity $A(w)=\spec_L(S)(w)$. A \emph{positive circulation of total $G$} is
a map $f:\supp\spec_L(S)\to\ints_{>0}$ balanced at every node with
$\sum_w f(w)=G$.
\end{definition}

The truth's spectrum $A$ is one such circulation, and the graph is strongly
connected because its edges are the edges of the single closed walk
$t\mapsto W_t(S)$.

\begin{lemma}[Same-length spelled candidates are positive circulations]
\label{lem:reduction}
Same-length Section~6.2 spelled candidates with observed support $V$ are exactly
the positive circulations of total $G$ on $X_S$.
\end{lemma}

\begin{proof}[Sketch]
A circular word $D$ of length $G$ has a balanced positive spectrum of total $G$;
if $\supp\spec_L(D)=V$ it is such a circulation. Conversely a positive
circulation on a strongly connected balanced multigraph makes it Eulerian, so an
Eulerian circuit spells a circular word $D$ of length $G$ with
$\spec_L(D)=f$. The rigidity theorem below proves uniqueness on the larger
circulation set, so it covers the spelled-circuit subclass a fortiori.
\end{proof}

\begin{theorem}[Short-repeat rigidity]
\label{thm:short-repeat}
If every length-$(L-1)$ window of $S$ occurs at most twice, then
$A=\spec_L(S)$ is the unique positive circulation of total $G$ on $X_S$.
\end{theorem}

\begin{proof}
Let $B$ be a positive circulation of total $G$ and $\delta=B-A$, a circulation
with $\sum_w\delta(w)=0$. Since $A(w)\le$ (occurrences of
$\operatorname{prefix}_{L-1}(w))\le 2$ and $B(w)\ge 1$, we have
$\delta(w)\ge-1$ on every edge.

Suppose $\delta\ne 0$. Then $N=\{w:\delta(w)<0\}$ is nonempty, and on $N$ we
have $\delta(w)=-1$, $A(w)=2$, $B(w)=1$. Take $w\in N$ with tail $v$. As
$A(w)=2$ and $\mathrm{out}_A(v)\le 2$, $w$ is the only outgoing edge at $v$, so
$\mathrm{out}_\delta(v)=-1$ and balance gives $\mathrm{in}_\delta(v)=-1$. Each
incoming edge has $\delta\ge-1$; if $v$ had two incoming edges they would each
have $A=1$ and hence $\delta\ge0$, contradicting $\mathrm{in}_\delta(v)=-1$. So
$v$ has a single incoming edge $e$ with $A(e)=2$, $\delta(e)=-1$, i.e.\ $e\in N$.
The symmetric argument at the head of $w$ shows every edge incident to a node
touched by $N$ lies in $N$; thus $N$ is a union of connected components of $X_S$
with no edge to $V\setminus N$. Strong connectivity forces $N=\emptyset$ or
$N=V$; the latter gives $\sum_w\delta(w)=-|V|<0$, contradicting
$\sum_w\delta(w)=0$. Hence $\delta\ge0$, and a nonnegative circulation with
$\sum_w\delta(w)=0$ is zero.
\end{proof}

The hypothesis of Theorem~\ref{thm:short-repeat} is exactly what the bridging
condition supplies. We record the two extension lemmas that transfer bridging
into short-repeat structure.

\begin{lemma}[Primitive extension]
\label{lem:primitive-extension}
Let $S$ be primitive and suppose some length-$(L-1)$ window occurs at least three
times. Then $S$ has a triple repeat of length at least $L-1$.
\end{lemma}

\begin{proof}[Sketch]
Choose three starts carrying the same length-$(L-1)$ window and extend the three
copies in both directions as long as all three remain equal. The process cannot
reach length $G$, since three equal length-$G$ windows would make $S$ invariant
under a nonzero shift, contradicting primitivity. At termination the preceding
symbols are not all equal and the following symbols are not all equal, which is
exactly the three-copy maximality condition; the length is at least $L-1$.
\end{proof}

\begin{lemma}[Periodic extension]
\label{lem:periodic-extension}
Let $S=P^k$ with minimal period $p<G$ and suppose some factor of $P$ of length at
least $L-1$ occurs at least twice in the circular word $P$. Then $S$ has a triple
repeat of length at least $L-1$.
\end{lemma}

\begin{proof}[Sketch]
A repeated factor of $P$ of length at least $L-1$ occurs at least four times in
$S$. Extending its two occurrence families maximally in both directions is
uniform by $p$-periodicity and cannot reach length $G$, because that would make
$r_2-r_1$ a period of $S$ and hence a smaller period of $P$. At termination the
three-copy maximality condition holds for any three occurrences meeting both
families.
\end{proof}

A length-$L$ read bridges a copy of a repeat of length $\ell$ only if
$\ell\le L-2$. Hence $R\in\Is$ forbids every triple repeat of length
$\ge L-1$. Combining this with the extension lemmas yields the main theorem.

\begin{theorem}[Rigidity under $\Is$]
\label{thm:rigidity}
If $S$ admits a read realization $R\in\Is$, then $A=\spec_L(S)$ is the unique
positive circulation of total $G$ on $X_S$.
\end{theorem}

\begin{proof}
If $S$ is primitive: $R\in\Is$ gives no triple repeat of length $\ge L-1$; by
Lemma~\ref{lem:primitive-extension} no length-$(L-1)$ window occurs three times;
Theorem~\ref{thm:short-repeat} applies. If $S=P^k$ with minimal period
$p<G$: by Lemma~\ref{lem:periodic-extension} every factor of $P$ of length
$\ge L-1$ occurs at most once, so the $p$ length-$L$ and length-$(L-1)$ windows
of $P$ are pairwise distinct and $X_S$ is a single directed $p$-cycle with
$A(w)=k$ on every edge. A positive circulation on a directed cycle is constant
by balance, and total $G=kp$ forces the value $k$, so $B=A$.
\end{proof}

\begin{corollary}[No strict same-length Section~6.2 counterexample]
\label{cor:same-length-zero}
Suppose $R\in\Is$ and the truth is Section~6.2-admissible, so
$\supp x=\supp\spec_L(S)$. Then every same-length Section~6.2 spelled candidate
$D$ satisfies $\spec_L(D)=\spec_L(S)$, and for every objective that is a function
of $(\spec_L(D),x)$,
\[
  \frac{\Lik(D;x)}{\Lik(S;x)}=1 .
\]
There is no strict same-length counterexample for any $G$, $L$, or alphabet.
\end{corollary}

\begin{proof}
By Theorem~\ref{thm:rigidity} the only positive circulation of total $G$ on
$X_S$ is $A$, so every same-length spelled candidate has spectrum $A$. Both
objectives depend on the candidate only through its spectrum, and the
observation is common, so the scores coincide.
\end{proof}

Only the triple-repeat clause of $\Is$ is used. The interleaved-repeat clause and
the coverage clause are not needed for this same-length conclusion. The equal
length assumption, by contrast, is essential and sharp; we return to it in
Section~\ref{sec:finite:length}.

\begin{proposition}[Beatability criterion]
\label{prop:beatability}
Fix $S,L$ and an observation $x$ with $\supp x=\supp\spec_L(S)$, and write
$A=\spec_L(S)$. If $A$ is the unique positive circulation of total $G$ on
$X_S$, the truth ties every same-length spelled candidate. If $B\ne A$ is
another positive circulation of total $G$, choose $w_0$ with $B(w_0)>A(w_0)$ and
put $x_M=x+M e_{w_0}$. Then the likelihood ratio equals
$\mathrm{ratio}_0\cdot(B(w_0)/A(w_0))^M$, which exceeds $1$ for all sufficiently
large $M$; the fixed-length binomial's domain is respected because $B(w_0)<G=N$
whenever $B$ is positive and non-constant.
\end{proposition}

Thus non-rigidity is exactly the same-length obstruction, and by
Theorem~\ref{thm:rigidity} the hypothesis $\Is$ removes it. The minimum
same-length non-rigid pair, found by exact search, is $S=AAAAB$,
$D=AABAB$ ($G=5$, $L=2$); its truth carries a long triple repeat and is
therefore not $\Is$-admissible. Skirting the triple-repeat clause is precisely
what permits the multiplicity reallocation.
\statusline{Theorem~\ref{thm:rigidity} and
Corollary~\ref{cor:same-length-zero} are \mathfact; the minimum non-rigid pair
and the bounded verifications are \computed, exhaustive in the cited scope.}

\subsection{Strict finite counterexamples}

The rigidity theorem is a positive same-length statement. Everything else in
this subsection is negative, and each witness is strict, hence robust to the
unresolved equivalence and tie conventions.

\begin{theorem}[Exact-multinomial counterexample over all circular candidates]
\label{thm:cex-exact}
Let $\Alphabet=\{A,B\}$, $L=3$, and let the truth be
\[
  S=AAABB \quad (G=5),
\]
with realized starts $0,1,4$, so the observed reads are $\{AAA,AAB,BAA\}$, each
once. Let
\[
  D=AAAAB \quad (|D|=5).
\]
Then $R\in\Is$ holds non-vacuously, and
\[
  \frac{\Lik_{\mathrm{ex}}(D;x)}{\Lik_{\mathrm{ex}}(S;x)} = 2 > 1 .
\]
Consequently the truth is not a maximizer for the exact multinomial over all
circular candidates, so both the maximizer and uniqueness schemas are false for
this objective.
\end{theorem}

\begin{proof}
The spectra on the observed types are $\spec_3(S)=(AAA:1,AAB:1,BAA:1)$ and
$\spec_3(D)=(AAA:2,AAB:1,BAA:1)$; $\spec_3$ has total $5$ for both. The
multinomial coefficient $\binom{3}{1,1,1}=6$ is common, and all candidates have
length $5$, so
\[
  \Lik_{\mathrm{ex}}(S;x)=6\cdot(1/5)^3=\frac{6}{125},\qquad
  \Lik_{\mathrm{ex}}(D;x)=6\cdot(2/5)(1/5)(1/5)=\frac{12}{125}.
\]
For $\Is$: the starts cover all five positions; the three $A$'s at positions
$0,1,2$ form a maximal length-$1$ triple repeat, bridged respectively by the
reads at starts $4,0,1$; and no two maximal repeat pairs have four alternating
starts, so the interleaved clause is vacuous. The all-bridged triple-repeat
clause is therefore genuinely exercised. The transfers to all circular
candidates by Lemma~\ref{lem:negative-transfer} because the witness is
same-length.
\statusline{\kernelcheck{} in the repository's
\texttt{FixedLengthExactCounterexample} module.}
\end{proof}

Repeating the start-$0$ read $k$ times gives the observed multiset
$\{AAA^k,AAB,BAA\}$ and ratio $2^k$, so the failure is a one-parameter family,
not a coincidental tie.

\begin{theorem}[Fixed-length binomial counterexample]
\label{thm:cex-binomial}
Let $\Alphabet=\{A,C\}$ inside the four-letter DNA alphabet, $L=3$, truth
$S=AAACC$ with realized starts $0,1,4$ (observed reads $AAA,AAC,CAA$), and
competitor $D=AAAAC$, both of length $5$. Under the literal Section~6.1 product
of binomial marginals with $N=5$, the zero-count factors retained,
\[
  \frac{\Lik_{\mathrm{bin}}(D;x)}{\Lik_{\mathrm{bin}}(S;x)}
  = \frac{1125}{512} > 1 .
\]
\end{theorem}

\begin{proof}[Sketch]
The instance is Theorem~\ref{thm:cex-exact} with $B$ relabeled to $C$, so the
same $\Is$ certificate transfers. The spectra are $S=(AAA:1,AAC:1,CAA:1)$ and
$D=(AAA:2,AAC:1,CAA:1)$ on the observed types. Carrying the factors for the
zero-multiplicity types changes the ratio from $2$ to $1125/512$:
\[
  \Lik_{\mathrm{bin}}(S;x)=\frac{452984832}{30517578125},\qquad
  \Lik_{\mathrm{bin}}(D;x)=\frac{7962624}{244140625}.
\]
\statusline{\kernelcheck{} in \texttt{FixedLengthBinomialCounterexample}. The
difference from the exact ratio shows that dropping or retaining the zero-count
factors changes the objective. By Lemma~\ref{lem:negative-transfer} the witness
refutes the maximizer statement over every circular candidate on which the
binomial objective is defined, since both $S$ and $D$ lie in the domain
($\spec_3(D)(AAA)=2<5$).}
\end{proof}

We now turn to the reverse-complement molecular reading of the Section~6.2 flow.
Here read types are molecule classes, the truth and competitor must both be
admissible bidirected circuits, and the objective is the literal Section~6.1
product with external $N$. The mechanism is different from the sequence-level
witnesses: reverse complementarity \emph{creates} multiplicity that the
observation does not force, and support equality then permits reallocation.

\begin{theorem}[Section~6.2 variable-length molecule counterexample]
\label{thm:cex-molecule}
Let $\Alphabet=\{A,T\}$ with involution $A\leftrightarrow T$, $L=3$, truth
$S=AAATT$ ($G=5$) with realized starts $0,1,4$, and competitor
$D=AAAATT$ ($|D|=6$). Then $R\in\Is$ holds; both $S$ and $D$ induce admissible
Section~6.2 bidirected circuits; and with external $N=5$,
\[
  \frac{\Lik_{6.1}(D)}{\Lik_{6.1}(S)} = \frac{9}{8} > 1 .
\]
\end{theorem}

\begin{proof}[Sketch]
Under $A\leftrightarrow T$ the molecule classes of $S$ are
$d_S=(AAA:1,AAT:2,TAA:2)$ and the observed classes are
$x=(AAA:1,AAT:1,TAA:1)$; for $D$ they are $d_D=(AAA:2,AAT:2,TAA:2)$. Both
spectra have the same support as $x$. The ratio of the $AAA$ factor is
$[(2/5)(3/5)^2]/[(1/5)(4/5)^2]=9/8$, and the other two classes have equal
multiplicities, so they contribute factor $1$. Bidirected-circuit feasibility and
the $\Is$ certificate are checked by exact computation and by an explicit
bidirected graph.
\statusline{\kernelcheck{} for the finite certificate and the strict inequality;
\computed{} for the bidirected graph certificate, which is checked by script and
not reproduced in Lean.}
\end{proof}

Allowing the lengths to be equal does not save the Section~6.2 reading. The
smallest same-length witness is $S=AAATAT$, $D=AAAAAT$ ($G=6$), for which the
exact ratio is $3$ and the binomial ratio is $5$, with both genomes
Section~6.2-feasible under the source's per-vertex lower bound.

\begin{remark}[A strengthening is open]
The witnesses above use the source's per-vertex lower bound, i.e.\ support
equality. Under the strictly stronger per-occurrence requirement
$d_D(w)\ge x(w)$ for every observed type, the same-length Section~6.2 statement
is \openstatus{} in the searched scope; the witnesses' truths are not candidates
under that strengthening. We flag this because it is the kind of
assumption-surface change that separates a refuted statement from an open one.
\end{remark}

\subsection{The candidate-length boundary}
\label{sec:finite:length}

The same-length restriction in the positive theorem is not cosmetic. On the next
slice, where $|D|>G$ is allowed, the truth need not be optimal even under $\Is$.

\begin{example}[Variable-length oriented boundary]
\label{ex:variable-length}
Take truth $S=AAATT$ ($G=5$), $L=3$, and competitor $D=AAAATT$ ($|D|=6$). The
supports agree and $\Is$ holds. With the observed counts
$x=\spec_3(S)+M e_{AAA}$,
\begin{center}
\begin{tabular}{@{}lll@{}}
\toprule
$M$ & exact multinomial & fixed-$N$ binomial ($N=5$) \\
\midrule
$0$ & $3125/3888<1$ & $81/128<1$ \\
$1$ & $15625/11664>1$ & $81/64>1$ \\
$2$ & $78125/34992>1$ & $81/32>1$ \\
\bottomrule
\end{tabular}
\end{center}
The boundary is exactly $n>G$: at $n=G$ the truth wins, and a single extra
observation of an over-represented type overturns it.
\statusline{\computed{}. The same word pair also appears as a
reverse-complement variable-length witness; here the orientation is fixed by
observing every truth type plus an extra $AAA$, so the mechanism is
observed-multiplicity skew rather than strand collapse.}
\end{example}

\subsection{Summary of the finite landscape}

\begin{table}[t]
\centering
\small
\begin{tabular}{@{}p{0.30\textwidth}p{0.24\textwidth}p{0.16\textwidth}p{0.20\textwidth}@{}}
\toprule
Statement & objective & candidate class & result \\
\midrule
same-length oriented Section~6.2, $R\in\Is$
  & exact / binomial
  & same-length spelled
  & \textbf{true} (Thm.~\ref{thm:rigidity}) \\
exact ML maximizer
  & exact multinomial
  & all circular
  & \textbf{false} (Thm.~\ref{thm:cex-exact}) \\
binomial ML maximizer
  & fixed-$N$ binomial
  & all circular (on its domain)
  & \textbf{false} (Thm.~\ref{thm:cex-binomial}) \\
Section~6.2 flow, molecule reads
  & Section~6.1 product
  & §6.2 circuits
  & \textbf{false} (Thm.~\ref{thm:cex-molecule}) \\
Section~6.2 flow, molecule, same length, per-vertex
  & exact / binomial
  & §6.2 spelled
  & \textbf{false} ($AAATAT\to AAAAAT$) \\
Section~6.2, per-occurrence strengthening
  & exact / binomial
  & §6.2 spelled
  & \openstatus \\
\bottomrule
\end{tabular}
\caption{The finite landscape. ``False'' means a strict counterexample over the
stated class; the evidence level (kernel-checked, exact computation, or script)
is recorded in the corresponding theorem or status line. By
Lemma~\ref{lem:negative-transfer} each same-length falsity transfers to all
circular candidates for that objective.}
\label{tab:finite}
\end{table}

Table~\ref{tab:finite} makes the shape of the answer visible. The 2016 question
is not one theorem awaiting one proof: it is a family of statements indexed by
objective, strand convention, candidate class, and length convention, some true,
some false, and at least one open. The negative results also expose the
\emph{mechanism}, which the next section isolates: the exact objective scores
empirical read-type frequencies, and a competitor can improve its score by
reallocating multiplicity toward over-represented observed types while
respecting the bridging constraint on the realized reads.

\section{Population maximum likelihood and a repaired uniqueness theorem}
\label{sec:population}

The finite failures of Section~\ref{sec:finite} have two different causes, and
separating them is what motivates the model of this section.

\emph{Structural non-identifiability} (Example~\ref{ex:ABAB}) is the failure to
determine the genome from the read distribution at all. Bridging conditions
address it, and Theorem~\ref{thm:rigidity} shows how far they go on the
same-length slice.

\emph{Finite-sampling fluctuation} is a different phenomenon. Even a candidate
that is structurally admissible can beat the truth because the realized sample
is unrepresentative. The next example isolates this mechanism in its purest
form.

\begin{example}[Finite frequencies can beat structural admissibility]
\label{ex:frequencies}
Let $S=AABBC$ ($G=5$) be a primitive circular genome and $L=3$. Realize the
starts $(0,3)$, giving the observed reads $\{AAB,BCA\}$. Let $D=AABC$ and note
$|D|=4$. Both $S$ and $D$ are primitive and satisfy P1 (no repeated
$(L-1)=2$-mer). The realized sample satisfies the truth-side
information-feasibility condition, yet under the candidate-intrinsic exact
multinomial objective
\[
  \frac{\Lik_{\mathrm{ex}}(D;x)}{\Lik_{\mathrm{ex}}(S;x)}
  = \frac{(1/4)^2}{(1/5)^2} = \frac{25}{16} > 1 .
\]
The wrong genome wins because the realized sample consists only of reads to
which the shorter candidate assigns probability $1/4$ while the truth assigns
$1/5$. Structural admissibility has not failed; the empirical frequencies have
fluctuated.
\statusline{\computed{}, exact rational arithmetic. This is the project's
finite witness for the sampling mechanism; it is not a counterexample to any
historical statement, since the candidate restrictions are project-level.}
\end{example}

The lesson is that strengthening the repeat axioms again would target the wrong
mechanism. The repair that removes empirical fluctuation is to replace the
finite sample by its population limit.

\subsection{Population likelihood}

\begin{definition}[Population read distribution and likelihood]
\label{def:population}
For a circular word $D$ and read length $L\le|D|$, the population read
distribution is
\[
  p_D(w) \;=\; \frac{\spec_L(D)(w)}{|D|},
\]
a probability distribution on length-$L$ words. For a fixed truth $S$, the
\emph{population log-likelihood} of a candidate $D$ is
\[
  \ellpop_S(D) \;=\; \sum_{w} p_S(w)\log p_D(w),
\]
with $\ellpop_S(D)=-\infty$ if some word with $p_S(w)>0$ has $p_D(w)=0$.
\statusline{\projchoice. This is the infinite-read (population) analogue of the
Medvedev--Brudno exact objective; it is a project-level formulation motivated by
the finite analysis, not an assumption of either historical paper.}
\end{definition}

\begin{lemma}[Gibbs/KL identity; the truth is always a population maximizer]
\label{lem:gibbs}
For any candidate $D$,
\[
  \ellpop_S(D)-\ellpop_S(S) = -\KL(p_S\,\|\,p_D) \;\le\; 0,
\]
with equality if and only if $p_D=p_S$. Consequently, over \emph{any} candidate
class containing $S$, the truth is a population maximum-likelihood genome, with
no repeat condition required.
\end{lemma}

\begin{proof}
$\ellpop_S(S)=\sum_w p_S(w)\log p_S(w)$, so
$\ellpop_S(D)-\ellpop_S(S)=\sum_w p_S(w)\log(p_D(w)/p_S(w))=-\KL(p_S\|p_D)$,
where the right side is $+\infty$ precisely when $\ellpop_S(D)=-\infty$. Gibbs'
inequality \cite{cover2006} gives $\KL\ge0$ with equality iff the distributions
agree. The last claim is immediate since $S$ is in the class.
\end{proof}

Lemma~\ref{lem:gibbs} is the conceptual heart of the repair. At the population
level the maximizer statement is free; the entire content of a uniqueness theorem
is then \emph{spectrum identifiability}: for which candidate classes is
$p_D=p_S$ equivalent to $D\approx S$?

\subsection{P2 plus primitivity remove spectrum scaling}

Population equality compares \emph{normalized} spectra. The project uses two
candidate-intrinsic admissibility predicates to cross the gap to ordinary
spectrum equality.

\begin{definition}[P1 and P2]
\label{def:P1P2}
Fix the read length $L$. An oriented circular word $D$ satisfies
\begin{itemize}[leftmargin=*]
  \item \textbf{P1} if no length-$(L-1)$ word occurs more than once in $D$;
  \item \textbf{P2} if $D$ has no triple repeat of length $\ge L-1$ and every
        interleaved maximal repeat pair has a constituent of length
        $\le L-2$.
\end{itemize}
P1 is stronger than P2 (no repeated $(L-1)$-mer implies the P2 repeat bounds).
Both are project-level, candidate-intrinsic conditions; neither is a historical
assumption.
\statusline{\projchoice. The threshold in P2 is chosen to match the
complete-spectrum uniqueness theorem at $K=L-1$ below.}
\end{definition}

\begin{lemma}[Primitive P2 spectra have gcd one]
\label{lem:scaling}
Let $S$ be a primitive P2-admissible circular word and $L\ge2$ with
$L\le|S|$. Let $c=\spec_L(S)$ viewed as an integer vector and let
$g=\gcd\{c(w):c(w)>0\}$. Then $g=1$.
\end{lemma}

\begin{proof}[Proof sketch]
Suppose $g>1$. Work in the order-$(L-1)$ de Bruijn graph: vertices are
$(L-1)$-mers and each occurrence of a length-$L$ word is a directed edge
with multiplicity $c(w)$. Because $S$ spells a closed trail using every
edge exactly once, this edge multiset is balanced (in-degree equals
out-degree at every vertex) and weakly connected on its support.
Dividing every multiplicity by $g$ gives $c'=c/g$, still integer-valued.
Balance is a homogeneous linear condition, so it is preserved by division;
the support is unchanged ($c'(w)>0\iff c(w)>0$), so weak connectivity on
the support is preserved as well. Hence $c'$ is a balanced connected edge
multiset and spells an Eulerian closed trail $W$ with
$\spec_L(W)=c'$ and $|W|=|S|/g$.
Repeating $W$ gives $W^g$ with $|W^g|=|S|$ and
$\spec_L(W^g)=g\cdot\spec_L(W)=c=\spec_L(S)$: each cyclic length-$L$
window of $W$ lifts to exactly $g$ windows of $W^g$.
Thus $S$ and $W^g$ share the same ordinary $L$-spectrum.
$S$ satisfies P2, hence Ukkonen's condition at $K=L-1$, so
Theorem~\ref{thm:BBT} applied to $S$ forces $S\sim W^g$ up to rotation.
But $W^g$ is a nontrivial whole-genome power while $S$ is primitive, and
primitivity is rotation-invariant --- a contradiction. So $g=1$.
Crucially, uniqueness is applied only to $S$; nothing requires $W^g$
itself to satisfy P2.
\end{proof}

\begin{corollary}[Normalized equality implies ordinary equality for primitive P2]
Let $S,D$ be primitive P2-admissible with $p_D=p_S$ for some $L\ge2$.
Then $\spec_L(D)=\spec_L(S)$.
\end{corollary}

\begin{proof}[Proof sketch]
$p_D=p_S$ rewrites as $|S|\,\spec_L(D)=|D|\,\spec_L(S)$, so the two
integer spectrum vectors are proportional. By the lemma both have gcd
$1$ over their positive entries, and proportional gcd-one vectors are
equal (indeed $|S|\mid|D|$ and $|D|\mid|S|$, so $|S|=|D|$).
\end{proof}

\begin{remark}[Why primitivity alone is not enough]
The unrestricted statement --- primitivity plus $p_D=p_S$ implies
$\spec_L(D)=\spec_L(S)$ --- is false. At $L=2$, the primitive words
$S=AAB$ ($|S|=3$) and $D=AAABAB$ ($|D|=6$) have
$\spec_2(S)=(AA{:}1,AB{:}1,BA{:}1)$ and
$\spec_2(D)=(AA{:}2,AB{:}2,BA{:}2)$, hence identical normalized spectra
$p_S=p_D=(1/3,1/3,1/3)$ but distinct ordinary spectra. The lemma above
uses P2 essentially, via the BBT uniqueness input.
\end{remark}

\subsection{The complete-spectrum uniqueness input}

The remaining step is a classical one, supplied by an external source.

\begin{theorem}[Bresler--Bresler--Tse complete-spectrum uniqueness]
\label{thm:BBT}
\srcfact{} \cite{bresler2013}. Construct the $K$-mer graph from the complete
$(K+1)$-spectrum of a circular genome. If the genome satisfies Ukkonen's
condition --- no triple repeat and no interleaved repeat pair of length at least
$K$ --- then the graph has a unique Eulerian cycle, which spells the genome up
to cyclic rotation.
\end{theorem}

Setting $K=L-1$ aligns the theorem's condition exactly with P2: maximal triple
repeats then have length at most $L-2$, and every interleaved maximal-repeat pair
has a constituent of length at most $L-2$.

\subsection{Population uniqueness for the oriented primitive P2 class}

\begin{theorem}[Population uniqueness, oriented primitive P2]
\label{thm:population}
Let $L\ge2$. Among primitive P2-admissible oriented circular candidates, the
true genome $S$ is the unique population maximum-likelihood genome up to cyclic
rotation.
\end{theorem}

\begin{proof}
Let $D$ be primitive and P2-admissible with $\ellpop_S(D)=\ellpop_S(S)$. By
Lemma~\ref{lem:gibbs}, $p_D=p_S$; by Lemma~\ref{lem:scaling} both spectra
have gcd one, hence the corollary gives
$\spec_L(D)=\spec_L(S)$. Both $S$ and $D$ satisfy P2, hence Ukkonen's condition
at $K=L-1$, so Theorem~\ref{thm:BBT} applied to the complete $(L)$-spectrum
gives a unique Eulerian cycle and therefore $D\sim S$. The truth is a maximizer
by Lemma~\ref{lem:gibbs}.
\end{proof}

\statusline{The project-side reduction (division/Eulerian gcd-one, issue
\#70) and the thin final tie-to-rotation wiring (issue \#73) are
\kernelcheck{} in \texttt{PopulationReduction} and
\texttt{PopulationUniqueness}, reusing the \#70 reduction with no
duplicated infrastructure. The Gibbs/KL tie characterization and the BBT
uniqueness step remain explicit external premises (\srcfact{}
\cite{cover2006}, \cite{bresler2013}). The gcd-one/division writeup is the
project's.}

\subsection{Scope and one failed proof route}

The theorem is stated for the \emph{oriented} spectrum model. It must not be
silently transferred to reverse-complement-collapsed molecule classes, whose
observation object is different: population equality there is a statement about
class-indexed distributions, and the identification of classes changes which
words are distinguishable.

A stronger statement was attempted directly within the project, via a matching
argument on the transition involution of the length-$(L-1)$ de Bruijn structure.
That route has a genuine gap: a crossing between two raw occurrence pairs need
not survive maximal extension as two interleaved \emph{maximal} repeats, and an
explicit primitive example with the same length-$3$ spectrum and no cyclic
relation exhibits the failure (it nevertheless fails P2 for a different
interleaved pair, so it refutes the proof step, not the theorem). The
Bresler--Bresler--Tse route above bypasses this gap; the failed route is recorded
because it shows that ``the same spectrum is unique'' is not automatic from
short local repeat structure.

Finally, the population result does \emph{not} answer the finite 2016 question.
It answers a different, cleaner question: once empirical frequencies are replaced
by the true read distribution, bridging-style repeat bounds plus primitivity do
recover oriented circular uniqueness. Whether the historical finite-sample
statement holds under any fixed source-supported reading remains governed by
Section~\ref{sec:finite}.

\section{Discussion: two failure modes, two repairs, and what remains open}
\label{sec:discussion}

\subsection{The shape of the answer}

The finite picture in Table~\ref{tab:finite} is neither a proof nor a refutation
of a single sentence. It is a map of an interface. Bridging conditions are a
\emph{structural} hypothesis; maximum likelihood is an \emph{optimization}
objective; the published sentence asks whether the first implies the second. Our
study shows that the answer depends on choices the sentence does not make, and
that the mathematically clean content lives in two distinct failure modes.

\paragraph{Structural/model failure.}
Without a repeat/read-length boundary the genome need not be identifiable at all
(Example~\ref{ex:ABAB}). Bridging conditions repair the structural ambiguity, and
Theorem~\ref{thm:rigidity} shows that on the same-length oriented Section~6.2
slice they are strong enough to make the truth's spectrum the only positive
circulation --- an exact rigidity, not just a reconstruction claim. Requiring
candidate length equal to the true length is, however, a strong extra axiom. It
is not supplied by Medvedev and Brudno's known-size parameter, and
Example~\ref{ex:variable-length} shows the conclusion fails as soon as it is
dropped.

\paragraph{Finite-sampling failure.}
Even structurally admissible genomes can lose to the truth's competitors because
empirical read frequencies fluctuate (Example~\ref{ex:frequencies}). This is a
different mechanism and needs a different repair. Population maximum likelihood
removes it at the maximizer level by construction
(Lemma~\ref{lem:gibbs}), and for the oriented primitive P2 class it also yields
uniqueness (Theorem~\ref{thm:population}).

The resulting conceptual progression is:

\begin{quote}
naive ML intuition
$\to$ structural repeat obstruction
$\to$ the literature's repeat/bridging boundary
$\to$ boundary-conditioned finite ML
$\to$ a positive same-length slice but negative variable-length neighbors
$\to$ replacing privileged true-length knowledge by intrinsic candidate checks
$\to$ finite failure surviving as sampling fluctuation
$\to$ population ML removing the fluctuation
$\to$ P2 plus primitivity restoring oriented circular uniqueness.
\end{quote}

This ordering matters: population data is not introduced to patch repeat
ambiguity, and intrinsic candidate checks are not introduced as assumptions of
the historical papers. They repair weaknesses exposed in sequence.

\subsection{What remains open}

\paragraph{Source questions.}
Which Medvedev--Brudno object the 2016 sentence denotes; which read-type
convention (and hence which equivalence) it intends; whether candidate length is
restricted; whether the conclusion is maximizer or uniqueness; and whether the
unretrieved publisher supplement contains a definition that narrows any of
these. These are not bookkeeping: each changes which formal statement the
published sentence asserts.

\paragraph{Mathematical questions.}
The same-length Section~6.2 statement under the per-occurrence strengthening
($d_D\ge x$) remains open in the searched scope. The single-strand reading of the
variable-length Section~6.2 statement remains open (bounded zero evidence). The
population uniqueness theorem is now supported by the P2-specific gcd-one
argument in Lemma~\ref{lem:scaling} together with the external
Bresler--Bresler--Tse uniqueness theorem; this chain remains source-supported
mathematics rather than a kernel-checked Lean theorem. The population model for
\emph{molecule} read types --- as opposed to oriented types --- is not addressed
here.

\paragraph{What would count as settlement.}
A positive settlement of the published question is a Lean proof of a formally
justified version of the published implication after an argument fixes the axes
of Table~\ref{tab:axes}. A negative settlement is a mathematically valid
counterexample satisfying the faithfully formalized bridging hypotheses while
violating the faithfully formalized likelihood conclusion. A computational
discovery is useful but not a settlement until the finite instance is
kernel-checked and the source correspondence is documented.

\subsection{Relation to the literature}

The closest published result is older than the 2016 question:
Bresler, Bresler, and Tse prove a \emph{necessary} condition --- if a
problematic interleaved pair or triple repeat has all copies unbridged, another
same-length genome has the same read likelihood. Shomorony et al.\ strengthen
the hypothesis to sufficient conditions for exact reconstruction and then ask
whether those conditions also imply ML optimality. A 2025 correctness-guarantee
paper still presents the information/correctness line and the
maximum-likelihood line as separate \cite{salmela2025}, which is meaningful
evidence that the 2016 question had not become a standard solved result. Our
negative witnesses are consistent with that picture: the implication is not
merely unproved but false under several precise readings, while the positive
same-length and population results show where the intuition does survive.

\subsection{Conclusion}

The 2016 question is best read not as one theorem awaiting one proof but as an
interface between a bridging hypothesis and a maximum-likelihood formulation
whose candidate, representation, length, and tie choices all matter. We have made
those choices explicit, proved a positive same-length rigidity theorem, refuted
several precise readings by strict kernel-checked witnesses, and presented a
repaired population model that restores oriented circular uniqueness under
primitivity and a complete-spectrum condition. The finite source-faithful
question remains distinct from the repaired result, and we have marked exactly
where the boundary lies.


\appendix
\section{Modeling choices explained}
\label{app:modeling}

This appendix explains why the model of Section~\ref{sec:model} has the shape it
does. It is written as a deliberate explanatory layer: it does not prove
anything, and each choice is labeled by its origin --- historical/source
mandated, mathematical idealization, operational approximation, or
project-level extension. Where an external treatment explains a convention
better than we can inline, we cite it rather than reproduce background.

\subsection{Why the genome is circular}

\paragraph{Source origin.}
Shomorony et al.\ work with a circular sequence ``for ease of exposition \dots
[to] avoid edge effects,'' so that every start position has a length-$L$ window
\cite[\S2]{shomorony2016}. Medvedev and Brudno's exact likelihood is likewise
stated for a ``circular genome'' \cite[\S6.1]{medvedev2009}.

\paragraph{What it buys.}
Circularity removes the two ends: the window map $t\mapsto W_t(D)$ is defined at
every position, the spectrum has total $|D|$ with no boundary correction, and the
window-support graph of Definition~\ref{def:graph} is balanced and strongly
connected because it is generated by a single closed walk. Several proofs would
need extra boundary terms without this.

\paragraph{What it approximates.}
A real chromosome is linear; the circular model is appropriate when the genome is
genuinely circular (many bacterial genomes) or when reads near the ends are rare
relative to the genome length, so edge effects are negligible. It is an
idealization, not a claim that genomes are circles.

\paragraph{Mandatory consequence.}
Cyclic-shift equivalence is forced by circularity (Remark~\ref{rem:cyclic}).
A formalization that models a circle but compares literal strings is internally
inconsistent.

\subsection{Why reads are fixed-length, error-free, and uniformly sampled}

\paragraph{Source origin.}
The 2016 exposition draws each of $N$ reads independently and uniformly from the
$G$ circular length-$L$ substrings of $s$ \cite[\S2]{shomorony2016}. Bresler et
al.\ use the same noiseless uniform model with bridging \cite{bresler2013}.

\paragraph{Why errors can be abstracted away in the source theory.}
Shomorony et al.\ deliberately separate an error-free information-theoretic phase
from sequencing-error handling; their algorithm uses exact overlaps, and their
double-strand implementation adds reverse-complement reads as preprocessing
\cite[\S4.1]{shomorony2016}. In an information-feasibility analysis this is a
modeling choice: it isolates the repeat-limited obstructions that determine the
information boundary. What the abstraction omits is real: sequencing errors
create spurious overlaps, turn unique repeats ambiguous, and interact with
coverage in ways a fixed $L$ cannot express. We therefore do not claim the
error-free model describes production assemblers.

\paragraph{Why uniform and independent.}
Uniformity over start positions makes the exact likelihood a multinomial with
probabilities $d_i/N(D)$; a non-uniform or position-dependent coverage profile
would destroy that conjugacy and change the objective. Independence makes the
joint likelihood a product, which is what enables both the exact multinomial and
the separable approximation. Neither is a biological law; both are the minimal
statistical assumptions under which the source's likelihood has the stated form.

\subsection{Why multiplicity is retained}

The likelihood is a function of the observed read-type \emph{counts}
$x(w)$, not of the set of distinct read strings. This is source-mandated: the
exact objective is the multinomial
$\frac{n!}{\prod_w x(w)!}\prod_w(d_w/N(D))^{x(w)}$. Collapsing reads to a set
would discard the frequencies that the objective scores and would make the
finite-sampling failure of Example~\ref{ex:frequencies} invisible. Keeping the
latent true placements separate from the observable counts is equally important:
bridging is a property of $(S,R)$, where $R$ records placements, while the
likelihood estimator may not see them.

\subsection{Oriented reads versus reverse-complement molecules}

\paragraph{Two genuinely different sources.}
Shomorony's theory uses oriented length-$L$ substrings; reverse complement enters
only as experimental preprocessing that \emph{adds} reads \cite[\S4.1]{shomorony2016}.
Medvedev and Brudno instead model a read as a DNA \emph{molecule}, an unordered
pair of reverse-complementary strands, represented once in the graph
\cite[\S3.1,\S4.1]{medvedev2009}. This is not notation: the two models have
different read-type spaces.

\paragraph{Consequences.}
Under oriented types, reverse-complementing a candidate relabels its type counts
by $t\mapsto\rc(t)$, and the candidate and its reverse complement need not tie.
In fact they tie only when the observed count vector is reverse-complement
symmetric, which fails in general, e.g.\ $D=AAAT$ at $L=1$ has $A$-count $3$
while $\rc(D)=TTTA$ has $A$-count $1$. Under molecule classes, a candidate and
its reverse complement \emph{always} tie because class counts are
reverse-complement invariant. Hence the read-type convention and the genome
equivalence are coupled (Section~\ref{sec:open}): molecule types force dihedral
equivalence, oriented types force only cyclic shift. The internal tension in
Medvedev and Brudno's Section~6.1 --- the text writes ``$4^k$ such variables''
while its vocabulary is $k$-molecules --- must be resolved per formalization
rather than glossed.

\subsection{Candidate length, known genome size, and same-length restriction}

These are three distinct assumptions, and conflating them is the most common
source-faithfulness error:

\begin{enumerate}[label=(\alph*),leftmargin=*]
  \item the likelihood may have a \emph{candidate-intrinsic} length $N(D)$
        (exact multinomial);
  \item it may be supplied an \emph{external} genome-size parameter $N$, used in
        the per-type binomial marginals (the Section~6.1 approximation);
  \item the \emph{candidate class} may be restricted to $|D|=G$.
\end{enumerate}
(a) and (b) are different objectives; (b) does not imply (c). Medvedev and
Brudno say ``we assume that the genome size is known''
\cite[\S6.1]{medvedev2009}, and a contemporaneous formulation fixes the assembly
length to the known genome length and justifies it by scale-degeneracy
\cite{ghodsi2016}. But the 2016 sentence itself fixes neither, so a same-length
theorem is a restricted result unless a source argument fixes the restriction.
The fixed-length witnesses of Section~\ref{sec:finite} are same-length precisely
to be robust: by Lemma~\ref{lem:negative-transfer} they refute the statement over
all circular candidates as well.

\subsection{Finite sample versus population}

A finite sample is what a sequencing experiment provides; a population is the
distribution it is drawn from. The two answer different questions. In the finite
regime the empirical frequencies fluctuate, and a wrong candidate can win
(Example~\ref{ex:frequencies}); the question is whether bridging is strong enough
to prevent this anyway, which Section~\ref{sec:finite} answers in the negative
for precise readings. In the population regime the fluctuation is gone, so the
maximizer statement is automatic and only identifiability remains
(Section~\ref{sec:population}). The population model is a project-level
idealization used to isolate the identifiability question; it is not the
historical finite-data model, and it does not settle it.

\subsection{What this paper deliberately does not model}

To keep the main narrative readable we excluded several directions that the
repository has not studied in depth: linear genomes with genuine end effects;
variable-length reads; noisy or adversarial reads beyond the source's
abstraction; and the Section~6.2 flow's non-contiguous assemblies as
sequence-valued candidates. We list them here rather than silently omitting
them, and we cite the adjacent literature on variable-length reads
\cite{hui2016} and safe/complete assembly \cite{tomescu2016,rahman2022} for
readers who want those directions.

\section{Epistemic ledger}
\label{app:ledger}

Every nontrivial claim in this paper belongs to one of a small number of
epistemic classes. This appendix consolidates them. The classes are:
\emph{source fact} (stated by a cited source), \emph{source-supported inference}
(a reading of a source that the source does not literally state),
\emph{project modeling choice} (introduced here, not in the literature),
\emph{mathematical result} (proved in the repository's notes),
\emph{kernel-checked} (checked in Lean),
\emph{verified computation} (exact-rational script), \emph{bounded computational
evidence} (search in a stated scope, not a proof), \emph{source gap} (the
sources do not determine it), and \emph{open}.

\subsection{Source facts and source gaps}

\begin{table}[h]
\centering
\small
\begin{tabular}{@{}p{0.46\textwidth}p{0.20\textwidth}p{0.26\textwidth}@{}}
\toprule
Claim & Status & Basis \\
\midrule
Circular oriented length-$L$ read model; independent uniform starts
  & source fact & \cite[\S2]{shomorony2016} \\
$\Is$ = coverage $\wedge$ all triple repeats all-bridged $\wedge$ interleaved
pairs bridged & source fact & \cite[Eq.~1]{shomorony2016} \\
Repeat/triple-repeat maximality; bridging strict on both sides
  & source fact & \cite{bresler2013} \\
Exact multinomial with candidate-intrinsic $N(D)$
  & source fact & \cite[\S6.1]{medvedev2009} \\
Fixed-$N$ binomial approximation; ``genome size known''
  & source fact & \cite[\S6.1]{medvedev2009} \\
Section~6.2 objects are bidirected flows, not necessarily sequences
  & source fact & \cite[\S6.2]{medvedev2009} \\
Reads are reverse-complement molecule pairs in MB09
  & source fact & \cite[\S3.1,\S4.1]{medvedev2009} \\
``$4^k$ variables'' vs.\ molecule vocabulary (internal tension)
  & source fact / tension & \cite[\S6.1]{medvedev2009} \\
Which MB09 object the 2016 sentence denotes
  & source gap & \cite[\S5]{shomorony2016} \\
Which read-type convention / equivalence; candidate length; maximizer vs.\
uniqueness & source gap & \cite[\S5]{shomorony2016} \\
Publisher supplement unretrieved & source gap & repository record \\
\bottomrule
\end{tabular}
\caption{Source facts and gaps. The gaps are what prevent any single theorem
from being advertised as settling the published sentence.}
\end{table}

\subsection{Mathematical and project results}

\begin{table}[h]
\centering
\small
\begin{tabular}{@{}p{0.48\textwidth}p{0.20\textwidth}p{0.24\textwidth}@{}}
\toprule
Claim & Status & Provenance \\
\midrule
Cyclic shift is forced for the strong schema; reverse complement forced iff
molecule types & mathematical fact & repository equivalence note \\
Negative results transfer to superclasses; positives do not
  & mathematical fact & Lemma~\ref{lem:negative-transfer} \\
Same-length oriented circulation rigidity from the no-long-triple-repeat
boundary, including the primitive/periodic split
  & kernel-checked project reduction &
  \texttt{OrientedRigidity}, \texttt{WordLayer}, \texttt{RepeatAdapter} \\
No strict same-length Section~6.2 counterexample, any $G,L,\Alphabet$
  & mathematical fact; final Lean wiring pending & Corollary~\ref{cor:same-length-zero}, issue \#74 \\
Uniqueness up to cyclic rotation on the truth-admissible same-length row
  & source-supported mathematical theorem &
  Corollary~\ref{cor:same-length-rotation}, \cite{bresler2013} \\
Beatability criterion; non-rigidity is the same-length obstruction
  & mathematical fact & Prop.~\ref{prop:beatability} \\
Gibbs/KL: truth is a population maximizer over any class
  & mathematical fact & Lemma~\ref{lem:gibbs} \\
Primitivity plus P2 gives gcd-one spectra (division + Eulerian + BBT),
  reducing normalized to ordinary spectra
  & kernel-checked project reduction &
  \texttt{PopulationReduction}, issue \#70 \\
Population uniqueness, oriented primitive P2
  & source-supported mathematical theorem; project-side reduction and
  final tie-to-rotation wiring kernel-checked &
  Thm.~\ref{thm:population}, \texttt{PopulationUniqueness} (issue \#73),
  \cite{bresler2013} \\
P1/P2 intrinsic admissibility; population objective
  & project modeling choice & this paper, Defs.~\ref{def:P1P2},
  \ref{def:population} \\
\bottomrule
\end{tabular}
\caption{Mathematical and project-level results. The same-length rigidity core
and its primitive/periodic repeat reduction are kernel-checked; the thin final
Section~6.2 theorem wiring is tracked separately. The population division/Eulerian
reduction remains in progress. External complete-spectrum uniqueness is not
claimed as kernel-checked.}
\end{table}

\subsection{Finite certificates and computations}

\begin{table}[h]
\centering
\small
\begin{tabular}{@{}p{0.44\textwidth}p{0.20\textwidth}p{0.28\textwidth}@{}}
\toprule
Claim & Status & Artifact \\
\midrule
$AAABB\to AAAAB$, exact ratio $2$, $\Is$ certificate
  & kernel-checked & \texttt{FixedLengthExactCounterexample} \\
$AAACC\to AAAAC$, binomial ratio $1125/512$
  & kernel-checked & \texttt{FixedLengthBinomialCounterexample} \\
$AAATT\to AAAATT$, §6.1 ratio $9/8$; $\Is$; §6.2 circuits
  & kernel-checked (+ script for the graph) &
  \texttt{Section62BridgingCounterexample} \\
$AAATAT\to AAAAAT$, same length, exact $3$, binomial $5$
  & kernel-checked (+ script) &
  \texttt{SameLengthSection62Counterexample} \\
$AAATT\to AAAATT$ variable-length oriented boundary
  & verified computation & rigidity script \\
Minimum same-length non-rigid pair $AAAAB/AABAB$
  & verified computation (bounded, exhaustive in scope) &
  rigidity script \\
$AABBC\to AABC$ sampling failure, ratio $25/16$
  & kernel-checked & \texttt{FiniteSamplingCounterexample} \\
Direct P2 matching proof's maximal-extension step is invalid
  & counterexample to a proof step, not the theorem & audit note \\
\bottomrule
\end{tabular}
\caption{Finite certificates. Kernel-checked statements concern the finite
instance and arithmetic; the source-correspondence claims are separate
source-supported inferences.}
\end{table}

\subsection{Unresolved register}

We reproduce the project's open items so that the boundary of this paper is
explicit: (1) the referent of the 2016 sentence; (2) the read-type convention and
hence the equivalence; (3) the candidate-length convention; (4) the
maximizer-vs-uniqueness reading; (5) the $4^k$ vs.\ molecule-class index count in
MB09 Section~6.1; (6) the unretrieved publisher supplement; (7) the
per-occurrence same-length Section~6.2 statement; (8) the single-strand
variable-length Section~6.2 statement; (9) the thin final population
tie-to-rotation wiring, kernel-checked as \texttt{PopulationUniqueness}
(issue \#73) with Gibbs/KL and BBT as explicit external premises; and
(10) the population model for molecule read types.

\section{Reproduction guide}
\label{app:reproduction}

This appendix records how to rebuild the paper and re-run every computational
artifact referenced above. It is written for a reader who wants to scrutinize
the results rather than take them on trust.

\subsection{Building this paper}

The manuscript is designed to build from a pinned, minimal TeX dependency set
without repository-specific absolute paths. From the repository root:

\begin{verbatim}
sh paper/build.sh
\end{verbatim}

The script changes into \texttt{paper/}, creates or reuses a local Guix profile
at \texttt{paper/.guix-profile} from \texttt{manifest.scm}, sources that
profile, and runs \texttt{latexmk} with \texttt{pdflatex} and \texttt{biber},
crossing references and the bibliography as needed. It writes
\texttt{paper/main.pdf} and the usual intermediate files, which are ignored by
version control.

As an alternative to the local-profile script, the manifest can be used with an
ephemeral Guix shell. From inside \texttt{paper/}, run:

\begin{verbatim}
guix shell -m manifest.scm -- \
  latexmk -pdf -interaction=nonstopmode -halt-on-error main.tex
\end{verbatim}

No absolute path appears in the sources.

\subsection{Repository verification artifacts}

The finite and computational claims correspond to artifacts in the repository:

\begin{itemize}[leftmargin=*]
  \item exact and binomial same-length witnesses:
        \texttt{scripts/} exact-rational verifiers and the Lean modules
        \texttt{AssemblyP1/FixedLengthExactCounterexample.lean} and
        \texttt{AssemblyP1/FixedLengthBinomialCounterexample.lean};
  \item Section~6.2 molecule witnesses:
        \texttt{scripts/verify\_se62\_bridging\_flow\_counterexample.py},
        \texttt{scripts/verify\_se62\_mb09\_bidirected\_graph.py},
        \texttt{scripts/verify\_samelength\_se62\_counterexample.py}, and the
        Lean modules \texttt{Section62BridgingCounterexample.lean} and
        \texttt{SameLengthSection62Counterexample.lean};
  \item same-length oriented rigidity:
        \texttt{scripts/verify\_oriented\_se62\_rigidity.py}, plus the
        kernel-checked abstract circulation core in
        \texttt{AssemblyP1/OrientedRigidity.lean} and circular-spectrum adapter
        in \texttt{AssemblyP1/WordLayer.lean}.  The repeat-theory bridge that
        discharges the primitive/periodic hypotheses remains separate, so these
        modules do not yet kernel-check the full $I_s$ rigidity theorem;
  \item equivalence coupling:
        \texttt{scripts/verify\_equivalence\_coupling.py}.
\end{itemize}

Every Python verifier is self-contained, uses exact \texttt{fractions.Fraction}
arithmetic, is deterministic, and exits non-zero on any failed assertion. The
Lean modules contain no \texttt{sorry}, \texttt{axiom}, or \texttt{admit};

\begin{verbatim}
lake build
\end{verbatim}

builds them under the pinned toolchain. Where a claim is kernel-checked, the
corresponding module is named in Appendix~\ref{app:ledger}; where a claim is
only computationally verified, it is labeled accordingly.

\subsection{Epistemic boundary of reproduction}

Reproduction verifies the \emph{finite mathematics}: spectra, likelihood ratios,
bridging certificates, and graph feasibility. It does not verify the
source-correspondence claims that identify a given finite objective or
candidate universe with the one the 2016 sentence intends, because those are
source-supported inferences rather than computations. The ledger in
Appendix~\ref{app:ledger} keeps the two apart.


\printbibliography

\end{document}
